% Separate post-CP37 candidate; B102's actual-coordinate lemma is an explicit dependency.
\subsection{A selected point field excludes all first-wild tame comparisons}

Throughout this subsection $p\ge5$, $q=p^2$, $N=q-1$, and
$K$ is complete discretely valued of mixed characteristic
$(0,p)$ with algebraically closed residue field,
$\mu_p\subset K$, and $e=e_K=p(p-1)$.
The groups are whole potentially supersingular height-two
Barsotti--Tate groups, and every torsion comparison is an actual
Galois-module isomorphism. Put $s=(p+1)e=pN$ and take the
common tame good field $L/K$ of degree $N$.

\begin{lemma}[A diagonal first coefficient acts on an actual point field]
\label{lem:B-RD103-point}
\status{Candidate; complete proof awaiting independent audits.}
Suppose an actual common $p^3$-torsion module is trivial modulo $p$
over $L$, and let $F$ be its full division field over $L$.
Write $W=\operatorname{Gal}(F/L)\subset1+pM_2(\mathbf Z/p^3)$
and let $W_1$ be its first matrix image.
Let $P$ be primitive, choose a basis with $P=e_1$ over
$\mathbf Z/p^3$, and suppose $W_1$ contains
$D=\operatorname{diag}(d_1,d_2)$ with $d_1\ne0$.
There is an actual $g\in W$ preserving $E=L(P)$, whose restriction
to $E$ has exact order $p^2$. More precisely,
\[
 \begin{gathered}
 g=1+pD+p^2C,\qquad
 gP=(1+pd_1)P+Q,\quad Q\in M[p](L),\\
 g^pP=P+p^2d_1P\ne P,\qquad g^{p^2}P=P .
 \end{gathered}                                      \tag{B-RD103.1}
\]
No normality of $E/L$ is assumed.
\end{lemma}
\begin{proof}
Lift $D$ to the actual group, with entries of the displayed
diagonal matrix represented in $\mathbf Z/p^3$. The remaining
second digit is some $C$. Thus $Q=p^2CP$ is a first-level point,
and all such points are $L$-rational.
Every $h\in\operatorname{Stab}_W(P)$ therefore fixes $gP$.
Consequently $L(gP)\subset L(P)$ inside the same actual $F$.
These two fields have the same degree over $L$, because $g$
is an automorphism fixing $L$. Hence $g(E)=E$; equivalently,
$g$ normalizes the actual point stabilizer.
The literal matrix identity
$(1+pD+p^2C)^p=1+p^2D\pmod{p^3}$, valid for odd $p$,
proves the remaining assertions since $d_1\ne0$.
Thus, if $S=E^{\langle g|_E\rangle}$, the actual extension
$E/S$ is totally ramified cyclic of degree $p^2$.
Both fields are complete and have the same algebraically
closed residue field.
\end{proof}

\begin{theorem}[First-wild recovery in the entire tame residual sector]
\label{thm:B-RD103-tame}
\status{Candidate; complete proof awaiting independent audits.}
Two groups under the standing hypotheses with isomorphic
actual $p^3$-torsion and tame reducible residual inertia have
equal finite special characters. No restriction on the order
of their full wild division image is imposed.
\end{theorem}
\begin{proof}
Suppose their finite characters differ. The square determinant
case is excluded by B77, and the complete arbitrary-image
Corollary~\ref{cor:B-RD84-first-wild} gives two-slope parameters
and exchanged high lines
\[
 B_i=qk_i,\quad1\le k_i\le p-2,\qquad
 H_1=\mathbf F_pe_1,\quad H_2=\mathbf F_pe_2 .
                                                        \tag{B-RD103.2}
\]
Both first-level modules are $L$-rational.
Let $c_i$ be the finite special tangent exponents.
Equality of the actual determinants gives equality of their
finite prime-to-$p$ determinants, so
$c_1\equiv c_2\pmod{p-1}$.
The line $H_1$ is low in model 2; its two actual tame characters
are $\zeta^{c_1-B_1}$ and $\zeta^{pc_2+B_2}$.
Their equality, reduced modulo $p-1$, gives
$k_1+k_2\equiv0\pmod{p-1}$.
The strict ranges in (B-RD103.2) then force
\[
                 k_1+k_2=p-1.                         \tag{B-RD103.3}
\]

By Lemma~\ref{lem:B-RD85-diagonal}, the first matrix layer has
a nonzero $11$ projection. The two residual tame characters
are distinct: equality would make their determinant a square,
contrary to the remaining determinant case. Averaging
conjugation over the finite tame cyclic group therefore
projects $W_1$ onto its diagonal part and kills both
off-diagonal characters. In particular there is an actual
first coefficient $D=\operatorname{diag}(d_1,d_2)\in W_1$
with $d_1\ne0$. This requires only the diagonal projector,
not separation of the two opposite off-diagonal weights.

Choose corresponding primitive third-level points $P_i$
such that $p^2P_i=e_1$ in their common first-level module.
Their stabilizers agree under the actual torsion isomorphism,
so they define one common field
$E=L(P_1)=L(P_2)$.
The first-level target is high in model 1 and low in model 2.
Write
\[
       \alpha_1=s-qk_1,\qquad\beta_2=pk_2 .
                                                        \tag{B-RD103.4}
\]
Both high inverse polygons are single $q$-segments: indeed
$B_1=qk_1>e$, equivalently
$p(p+1)k_1>N$, puts the degree-$p$ vertex strictly above
the high-target chord; the next target is smaller.
The low inverse polygons are single $q$-segments as well.
The full integral coefficient proof is the actual inverse
polygon computation in B69 and B85.
Thus every such selected point has
\[
 v_L(T_1(P_1))=\alpha_1/q^2,\qquad
 v_L(T_2(P_2))=\beta_2/q^2=k_2/p^3 .
                                                        \tag{B-RD103.5}
\]
Since $p\nmid k_2$, its reduced valuation denominator is
$p^3$, forcing $p^3\mid[E:L]$.
On the other hand $p^2P_i$ is $L$-rational and its inverse
fiber under $[p^2]$ has $p^4$ elements. The actual point
orbit therefore has size at most $p^4$.
Moreover $[E:L]$ is a power of $p$, because it is a
subfield degree in the actual $p$-group division extension
$F/L$. Hence
\[
                         [E:L]\in\{p^3,p^4\}.           \tag{B-RD103.6}
\]
This is a statement about a selected point field, not the
size of the full division group.

The preceding lemma applies after lifting the residual
basis with $P_i$ as its first vector. It supplies one
actual order-$p^2$ automorphism $g$ of $E$, compatible
with both coordinates. The point $(g-1)P_i$ has order
$p^2$, and its $p$-multiple is $d_1p^2P_i$.
This is high in model 1 and low in model 2.
The same inverse polygon gives its $L$-valuation as
$\alpha_1/q$ and $\beta_2/q$, respectively.
The displacement by $g^p$ is the nonzero first-level
point $d_1p^2P_i$.
If $A_i=v_E(T_i(P_i))$, it follows that
\[
       v_E(gT_i(P_i)-T_i(P_i))=qA_i,\qquad
       v_E(g^pT_i(P_i)-T_i(P_i))=q^2A_i .
                                                        \tag{B-RD103.7}
\]
Formal subtraction and ordinary coordinate subtraction
differ by an integral unit, so these are the displayed
field-coordinate valuations.

If $[E:L]=p^3$, then
$A_1=N-pk_1$ and $A_2=k_2$, both prime to $p$.
They are different, since the first is a high-target
value and in fact
$A_1-pA_2=p-1>0$ by (B-RD103.3).
The exact prime-valuation equality of
Lemma~\ref{lem:B-RD60-break}, applied to this same $g$
in $E/S$, would give
$b(g)=NA_1=NA_2$. This is impossible.

If $[E:L]=p^4$, then $A_1=\alpha_1$ and $A_2=\beta_2$.
They have exact $p$-adic order one, and (B-RD103.3) gives
\[
                  A_1-pA_2=e>0 .
                                                        \tag{B-RD103.8}
\]
Moreover $q^2A_1<v_E(p)=p^4eN$, since $\alpha_1<eN$.
All hypotheses of the actual-coordinate
Lemma~\ref{lem:B-RD102-scalar} now hold for $E/S$ and $g$.
Its contradiction excludes this remaining field degree.
This proves the theorem.
\end{proof}

\begin{remark}[The remaining first-wild gate]
\label{rem:B-RD103-scope}
\status{Candidate; no nontrivial residual-wild sector is settled here.}
The argument treats all primes $p\ge5$ and both possible
opposite-weight patterns. Its only use of a cyclic group
is the actual order-$p^2$ subgroup acting on the selected
point field. Neither that point field nor the full division
field is assumed to be an abelian extension of $L$.
In particular no finite Cartan image is promoted to a
characteristic-zero CM action.

Together with the proved B98 irreducible recovery, the
standing dimension and mixed exclusions of B76, and the
earlier ordinary/filtered reductions, the only possible
first-wild third-level obstruction at a prime beyond the
already settled $p=5$ has nontrivial residual wild inertia
in a whole potentially supersingular comparison.
No such pair is constructed or excluded by this theorem.
The actual-coordinate lemma B102 is an explicit candidate
dependency; both proofs require complete independent audits
before promotion. A finite supplement will be retained
separately and will not replace any field or Newton argument.
\end{remark}
